% Taxer mieux plutôt que davantage : quatre réformes à recette constante (rapport Sorbon, partie III)
% Note unique issue de la fusion des notes ir_csg, foncier, is et pfu (novembre 2026).
% Compilation depuis fiscalite/synthese/ : pdflatex note-reforme-fiscale.tex (deux passes)
% Chiffres : sorties/perte_seche_ircsg_purs.txt (IR-CSG), ../foncier/sorties/simulation.txt,
% ../is/sorties/simulation_is.txt, sorties/perte_seche_autres.txt ; méthode : ../revue/methode_perte_seche.md
\documentclass[11pt,a4paper]{article}
% Préambule commun au brief et à la revue de littérature sur la taxe sur la valeur foncière.
% Même composition que le cours (cours/refaire-l-impot.tex) ; figures « nues » produites par
%   FIG_NU=1 python lvt/graphiques_lvt.py
\usepackage[T1]{fontenc}
\usepackage[utf8]{inputenc}
\usepackage{lmodern}
\usepackage[french]{babel}
\usepackage[a4paper,left=2.8cm,right=2.8cm,top=2.6cm,bottom=2.8cm]{geometry}
\usepackage{graphicx}
\usepackage[export]{adjustbox}
\usepackage{booktabs}
\usepackage{tabularx}
\usepackage{array}
\usepackage{xcolor}
\usepackage{titlesec}
\usepackage[labelfont=bf,labelsep=period,font=small,singlelinecheck=false,
            justification=raggedright,skip=6pt]{caption}
\usepackage{fancyhdr}
\usepackage{enumitem}
\usepackage{microtype}
\usepackage{textcomp}
\usepackage{amsmath}
\providecommand{\euro}{\texteuro}
\usepackage[hidelinks]{hyperref}

\definecolor{encre}{HTML}{1A1A19}
\definecolor{bleu}{HTML}{003399}
\definecolor{vert}{HTML}{1F7A4D}
\definecolor{ocre}{HTML}{C79100}
\definecolor{brique}{HTML}{993333}
\definecolor{filet}{HTML}{C8C6BE}

\titleformat{\section}{\normalfont\large\bfseries\color{bleu}}{\thesection.}{0.6em}{}
\titleformat{\subsection}{\normalfont\normalsize\bfseries\color{bleu}}{\thesubsection}{0.6em}{}
\titlespacing*{\section}{0pt}{2.4ex plus 1ex minus .2ex}{1.2ex plus .2ex}
\titlespacing*{\subsection}{0pt}{2ex plus 1ex minus .2ex}{0.8ex plus .2ex}

\pagestyle{fancy}
\fancyhf{}
\renewcommand{\headrulewidth}{0.4pt}
\fancyhead[R]{\small\thepage}

\setlength{\parindent}{0pt}
\setlength{\parskip}{0.62em}
\linespread{1.05}
\renewcommand{\arraystretch}{1.2}
\setlist{itemsep=2pt,topsep=2pt,parsep=0pt,leftmargin=1.4em}

\newenvironment{encadre}[1]
  {\par\medskip\noindent\begin{minipage}{\linewidth}%
   \color{encre}\setlength{\parskip}{0.5em}%
   \hrule height 0.8pt width \linewidth \vspace{6pt}%
   {\footnotesize\bfseries\color{vert}\MakeUppercase{#1}}\par\vspace{-2pt}%
   \small}
  {\par\vspace{4pt}\hrule height 0.4pt width \linewidth\end{minipage}\par\medskip}

\newcommand{\figbloc}[3]{%
  \begin{figure}[htbp]\centering
  \caption{#2}\label{#3}
  \includegraphics[max width=\linewidth,max height=0.40\textheight]{../figures/nu/#1.pdf}
  \IfFileExists{../figures/nu/#1_note.tex}{\par\vspace{4pt}%
    \parbox{\linewidth}{\raggedright\scriptsize\color{encre}\setlength{\parskip}{2pt}%
    \input{../figures/nu/#1_note.tex}}}{}%
  \end{figure}}

\newcommand{\pc}{\,\%}

\renewcommand{\figbloc}[3]{%
  \begin{figure}[!htbp]\centering
  \caption{#2}\label{#3}
  \includegraphics[max width=\linewidth,max height=0.28\textheight]{../figures/nu/#1.pdf}
  \par\vspace{3pt}\parbox{\linewidth}{\raggedright\scriptsize\color{encre}\setlength{\parskip}{2pt}%
    \input{../figures/nu/#1_note.tex}}%
  \end{figure}}
\newcommand{\figgrand}[3]{\figbloc{#1}{#2}{#3}}
\usepackage{needspace}
\usepackage{etoolbox}
\preto\section{\needspace{8\baselineskip}}
\preto\subsection{\needspace{5\baselineskip}}
\renewcommand{\topfraction}{0.9}\renewcommand{\bottomfraction}{0.8}
\renewcommand{\textfraction}{0.08}\renewcommand{\floatpagefraction}{0.85}
\setcounter{topnumber}{3}\setcounter{totalnumber}{5}
\fancyhead[L]{\small\itshape Taxer mieux plutôt que davantage}

\begin{document}
\renewcommand{\figurename}{Figure}\renewcommand{\tablename}{Tableau}

\thispagestyle{empty}
\vspace*{0.6cm}
{\footnotesize\color{ocre}\bfseries NOTE DE TRAVAIL\quad·\quad NOVEMBRE 2026\par}
\vspace{8pt}
{\fontsize{22}{26}\selectfont\bfseries\color{bleu}Taxer mieux plutôt que davantage\par}
\vspace{6pt}
{\large Quatre réformes à recette constante : impôts sur le revenu, immobilier, sociétés, revenus du capital\par}
\vspace{8pt}
{\small Colin Baget. Droit de 2026 ; données fiscales de la DGFiP (revenus 2024, fiscalité locale 2024, IS
2022), ventes immobilières DVF 2023-2024, comptes nationaux de l'Insee ; programmes \texttt{fiscalite/synthese/}
et dossiers \texttt{ir\_csg/}, \texttt{foncier/}, \texttt{is/}, \texttt{pfu/}. Vérifications et revues :
\texttt{fiscalite/revue/}.\par}
\vspace{6pt}
{\color{bleu}\hrule height 1.2pt}
\vspace{4pt}

\section*{Introduction}

Cette note examine quatre réformes avec un seul critère : à recette à peu près constante, quel système
fiscal gêne le moins la croissance, c'est-à-dire décourage le moins de travailler, d'investir, de déménager
ou de déclarer ses revenus ? La redistribution n'est pas l'objectif, mais deux contraintes s'imposent. La
somme des réformes ne doit pas alourdir nettement la facture des plus grandes fortunes : leur départ coûterait
plus qu'il ne rapporterait. Et le système doit être, à la marge, plus favorable aux familles avec enfants.
Il s'agit de taxer mieux, pas davantage.

La méthode est la même pour les quatre impôts. On part d'une référence simple : l'assiette la plus large,
le même taux pour tout le monde, sans niche. Puis on fait apparaître chacun des endroits où un choix doit
être fait (une franchise pour les revenus d'activité, un traitement des familles, la vitesse de la
déduction de l'investissement, le sort des plus-values logées dans une holding) et on chiffre chaque
choix de la même façon : ce qu'il coûte en recette, ce qu'il coûte ou rapporte en perte sèche, et qui
paie. La note ne tranche pas ces arbitrages à la place du lecteur ; elle dit ce que chacun coûte.

\begin{encadre}{Comment la perte sèche est calculée}
\textbf{Ce qu'elle mesure.} Un impôt fait renoncer à des choses qui auraient valu plus qu'elles ne coûtaient :
une heure de travail, une machine, une vente de logement. La perte sèche est la valeur de ces renoncements,
au-delà de l'impôt lui-même. Elle se mesure par la réaction des quantités taxées à l'impôt.

\textbf{La formule.} La perte sèche ordinaire, $\tfrac12\,\varepsilon\,\tau^2/(1-\tau)$ fois la base, suppose une seule marge et aucun autre impôt
dessus. On utilise sa généralisation (Harberger ; Hines, 1999 ; Chetty, 2009) : la variation de perte sèche
est la somme, sur chaque marge touchée, du coin fiscal total qui frappe cette marge multiplié par la
variation de la quantité taxée. Un euro de salaire perdu coûte à la collectivité tout ce qu'il aurait payé :
impôt sur le revenu et CSG, mais aussi cotisations non contributives et prestations perdues. Trois
conséquences comptent :
\begin{itemize}\setlength{\itemsep}{0pt}
\item \emph{le coin total} : relever l'impôt là où les cotisations et le retrait des prestations prélèvent
déjà 60~\% d'un euro de salaire en plus (autour de 1 à 1,5 Smic) coûte beaucoup plus que le calcul sur le
seul impôt ne le dit ;
\item \emph{les déplacements entre impôts} : un revenu qui quitte le barème pour l'impôt sur les sociétés ou le
prélèvement forfaitaire n'est perdu que pour l'écart de taux ; une réaction de calendrier (distribuer un
dividende un an plus tard) ne coûte presque rien ;
\item \emph{la décision de travailler} : on compte aussi la réaction de ceux qui hésitent à prendre un emploi,
qui dépend de ce que l'impôt prélève sur l'ensemble du premier salaire.
\end{itemize}

\textbf{Comparer des variantes.} Quand deux variantes ne lèvent pas la même recette, la différence doit être
levée ailleurs. On la valorise au « prix d'un milliard » levé par le barème fusionné : 0,16 à 0,60~Md€ de perte
sèche par milliard selon le coin retenu, 0,38 au centre (méthode et paramètres : \texttt{revue/methode\_perte\_seche.md}).

\textbf{Ce qui reste fragile.} Les élasticités, et surtout deux paramètres sans estimation française : la part
de la réaction des hauts revenus qui n'est qu'un déplacement vers l'impôt sur les sociétés, et la part de
l'ajustement des distributions de dividendes qui n'est qu'un report. Chaque résultat est donné avec sa
fourchette.
\end{encadre}


\section{Fusionner l'impôt sur le revenu, la CSG et la CRDS}\label{sec:ircsg}
% Partie IR-CSG PROVISOIRE (copie de ../ir_csg/note-ir-csg.tex), à réécrire
\begin{encadre}{Partie provisoire : chiffres en cours de recalcul}
Cette partie reprend la note sur la fusion telle qu'elle était avant la vérification de ses calculs. Celle-ci
(\texttt{revue/verif\_ir\_csg.md}) a trouvé un défaut de fond : chaque tranche de revenu est représentée par un
foyer qui touche à la fois des pensions et des revenus d'activité, si bien que la franchise d'activité profite
aussi aux retraités. Conséquences connues : les pensions paient au moins 28~Md€ de plus, et non 13 ; avec des
foyers séparés selon la source de revenu, le même barème rapporterait jusqu'à 243~Md€ au lieu de 220 ; les
parents isolés sont deux fois trop nombreux (4,0 millions au lieu de 2,0) et le quotient familial des enfants
coûte environ 12,9~Md€, non 13,9. Le modèle est en cours de reconstruction avec des foyers d'actifs et de
retraités séparés ; la partie sera réécrite à partir d'une référence à assiette large et même barème pour
tous, chaque arbitrage (franchise, pensions, famille, imposition commune, taux supérieur) étant chiffré avec
la perte sèche généralisée de l'encadré de méthode. Les chiffres ci-dessous ne doivent pas être repris.
\end{encadre}

La France dispose aujourd'hui de trois grands impôts sur le revenu : l'impôt sur le revenu (IR,
92~Md€ pour les revenus 2024), la CSG (153~Md€) et la CRDS (9~Md€). Le premier est fortement progressif
mais mité par une multitude de niches, abattements et dispositifs particuliers. La deuxième ressemble
davantage à l'idée qu'un économiste débarquant de Mars se ferait d'un impôt efficace : une assiette
large, des taux relativement faibles et peu de niches. La troisième obéit à une logique assez proche.

À première vue, rien de bien compliqué pour rationaliser l'ensemble : on fusionne les trois impôts, on
définit un nouveau barème qui rapporte autant, et l'on dispose d'un impôt unique, plus simple et plus
transparent. Mais si l'on se contente de conserver toutes les règles existantes sous une appellation
commune, on n'a pas simplifié grand-chose. \textbf{Le véritable intérêt de la réforme suppose d'harmoniser
les assiettes, les exemptions et les modalités de calcul. Et c'est là que les difficultés commencent.}

Car la plupart des exceptions fiscales, aussi critiquables soient-elles, ont une justification
économique ou sociale, au moins à l'origine. Prenons l'abattement de 10~\% sur les pensions de retraite.
Dans le projet de loi de finances pour 1978, le gouvernement proposait un abattement de 5~000~F accordé une
seule fois, au départ en retraite : les nouveaux retraités devaient acquitter un impôt calculé sur leurs
revenus d'activité, alors même que leurs ressources venaient de diminuer. L'Assemblée nationale a
transformé cet allégement ponctuel en une déduction permanente de 10~\% sur toutes les pensions (Sénat,
1977). Près de cinquante ans plus tard, le prélèvement à la source a fait disparaître le problème initial,
mais l'avantage demeure et coûte 5,3~Md€ par an ; le projet de loi de finances pour 2027 proposait seulement
d'en abaisser le plafond à 3~000~€ par foyer, article supprimé par l'Assemblée en octobre 2026. Le crédit d'impôt pour l'emploi à domicile, créé en 1991
pour créer des emplois et faire déclarer un travail souvent dissimulé, coûte aujourd'hui 7~Md€. Chacun de
ces avantages a ses raisons d'être, mais leur accumulation finit par coûter cher, et certaines
justifications ont perdu de leur pertinence.

\textbf{Le quotient familial pose une question plus fondamentale encore.} On peut le considérer comme une
aide à la natalité, auquel cas il est difficile de justifier qu'il bénéficie davantage aux familles
aisées. Mais on peut aussi y voir un principe de justice fiscale : selon les conventions de l'Insee, un
couple avec trois jeunes enfants gagnant 10~000~€ par mois a un niveau de vie comparable à celui d'un
couple sans enfant gagnant 6~250~€. Supprimer purement et simplement le quotient familial reviendrait à
opérer une redistribution considérable des familles avec enfants vers les contribuables sans enfant.
Même difficulté avec la CSG, dont les taux varient selon la nature des revenus et le statut des
contribuables : pourquoi les retraités paieraient-ils moins que les actifs à revenus comparables ? On
peut défendre cette différence, mais elle ne relève pas d'une nécessité économique. La supprimer revient
à transférer une partie de la charge des actifs vers les retraités, et c'est ce transfert qui, on le
verra, finance l'essentiel de la réforme.

Nous avançons par étapes, à rendement inchangé, chacune corrigeant le principal défaut de la précédente.
La première harmonise brutalement les règles en gardant les barèmes actuels. La deuxième remplace leur
empilement par un barème unique sans marche, le plus proche possible de ce que chacun paie aujourd'hui. La
troisième cherche le barème qui décourage le moins le travail et la déclaration des revenus, sans autre
considération. Notre proposition corrige enfin ce que ce barème a d'inacceptable : une forte hausse pour les
foyers modestes, et un gain bien plus grand pour une personne seule aisée que pour une famille de même
revenu. \textbf{Reste la famille} : l'IR repose sur le foyer, la CSG et la CRDS sont individuelles. Nous
comparons quatre réponses, de l'absence de toute aide familiale dans l'impôt à un quotient familial
appliqué à l'impôt fusionné. \textbf{Car derrière une réforme fiscale apparemment technique se cache
finalement une question assez simple : quelles différences entre contribuables l'État considère-t-il comme
suffisamment légitimes pour justifier qu'ils ne paient pas le même impôt ?}

\begin{encadre}{Comment les étapes sont chiffrées}
\textbf{Ce qui est remplacé.} L'IR, la CSG et la CRDS sur les revenus imposés au barème de l'IR : salaires,
pensions, bénéfices des indépendants, loyers, allocations chômage. Ils rapportent 220~Md€ : 80~Md€ d'IR
(une fois retiré celui qui porte sur les revenus du capital au prélèvement forfaitaire), 132~Md€ de CSG et
7~Md€ de CRDS. Chaque étape rapporte exactement autant. Les revenus du capital au prélèvement forfaitaire
unique gardent leur régime.

\textbf{L'assiette commune} est le revenu après cotisations sociales, le même pour un salarié, un
indépendant, un retraité ou un chômeur, sans réduction, abattement ni taux réduit. Les niches supprimées
représentent au moins 24~Md€ (annexe) ; la charge actuelle du modèle est l'impôt réellement payé, niches
déduites. L'impôt est individuel : chaque adulte est imposé sur son propre revenu.

\textbf{Les données.} Les 25 tranches de revenu fiscal de l'IRCOM (DGFiP, revenus 2024), réparties en
personnes seules avec ou sans enfant et couples sans enfant ou avec deux enfants, de façon à retrouver les
18~millions d'enfants à charge ; partage des revenus dans les couples selon l'Insee (Morin, 2014). Les foyers
types sont calculés avec OpenFisca-France ; sans niche, ils paient souvent plus aujourd'hui qu'un foyer réel.

\textbf{La perte sèche} est le revenu que les contribuables cessent de gagner ou de déclarer à cause de
l'impôt, net de ce qu'ils y gagnent en temps libre ; elle croît avec le carré du taux marginal (Feldstein,
1999). Elle porte sur les seuls revenus du travail, avec des élasticités de 0,15 pour la moitié la moins
aisée à 0,40 pour le 1~\% le plus aisé : environ 18~Md€ par an aujourd'hui, 8 à 31~Md€ selon les
élasticités.

\textbf{Les perdants} sont les foyers dont la charge augmente de plus de 500~€ par an. Dans une tranche,
tous les foyers d'un même type ont ici la même composition de revenus : le décompte sous-estime la
dispersion réelle et ne vaut que comme ordre de grandeur.
\end{encadre}

\subsection{Aujourd'hui : trois barèmes empilés}

Pour une personne seule salariée, l'IR, la CSG et la CRDS prélèvent 9,5~\% du salaire brut jusqu'à
environ 23~000~€, puis le taux marginal saute à 21~\%, 32~\% dès 40~000~€, 40~\% au-delà de 110~000~€
(figure~\ref{fig:auj}). Pour un couple avec deux enfants, le quotient familial et la décote maintiennent le
taux marginal à 9,5~\% jusqu'à 60~000~€. Les retraités subissent des sauts sans justification : quand leur
revenu fiscal franchit un seuil, le nouveau taux de CSG s'applique à toute leur pension, et 1~000~€ de pension
en plus coûtent 620 à 770~€.

\figbloc{ircsg_v01_aujourdhui}{Taux moyen et taux marginal aujourd'hui, trois foyers types}{fig:auj}

\subsection{Étape 1 : harmoniser sans toucher aux barèmes}

\textbf{La règle.} On garde le barème de l'IR (0, 11, 30, 41 et 45~\% avec les seuils des revenus 2024) et on
l'applique à chaque adulte, sur l'assiette commune, après une franchise de 6~000~€ par enfant déduite du
revenu du parent qui gagne le plus. On garde aussi un prélèvement proportionnel sur tout le revenu,
l'héritier de la CSG et de la CRDS, dont le taux est calé pour obtenir la même recette.

\textbf{Le taux proportionnel tombe à 3,0~\%.} Appliqué à une assiette sans niche, sans abattement de
10~\%, sans quotient ni décote, le barème de l'IR rapporte à lui seul 179~Md€, plus du double des 80~Md€
d'aujourd'hui. Il ne reste que 41~Md€ à trouver avec le prélèvement proportionnel, au lieu de 140~Md€ de CSG
et de CRDS. L'harmonisation à barèmes constants rend donc l'ensemble \emph{beaucoup plus} progressif : son
effet redistributif dépasse d'un quart celui d'aujourd'hui.

\figgrand{ircsg_x01_harmonisation}{Étape 1 : taux moyen et taux marginal}{fig:ex1}

\textbf{Qui gagne, qui perd.} Les neuf premiers déciles paient moins, surtout les déciles 5 à 7 (790 à
870~€ de moins par an), parce que le prélèvement proportionnel, qui frappe dès le premier euro, baisse
fortement. Le dernier décile paie 3~040~€ de plus en moyenne (2,4~points de revenu). Les pensions paient
8,8~Md€ de plus, les allocations chômage 3~Md€ ; un retraité seul à 17~000~€ de pension, aujourd'hui exonéré,
paierait 1~110~€. Les couples à un seul revenu paient beaucoup plus : avec deux enfants et 80~000~€, un
couple gagne 4~130~€ s'il partage ses revenus à 60/40, mais paie 2~590~€ de plus si un seul conjoint les
gagne, car l'imposition individuelle supprime le quotient conjugal. Au total, 14~\% des foyers perdent plus
de 500~€ par an.

\textbf{Le défaut : on a supprimé les niches, mais gardé les marches} (figure~\ref{fig:ex1}). Le taux
marginal d'une personne seule saute de 3~\% à 12~\% vers 13~000~€ brut, puis à 29~\% vers 33~000~€ et à
40~\% vers 95~000~€. La perte sèche reste celle d'aujourd'hui, 18~Md€.

\subsection{Étape 2 : un barème unique sans marche}

\textbf{La règle.} Le taux marginal monte en ligne droite entre quelques seuils, sans aucune marche, jusqu'à
50~\% pour les plus hauts revenus. Parmi les barèmes de ce type, on retient celui qui s'écarte le moins de
ce que chaque foyer paie aujourd'hui, à même recette. Il prélève, par adulte et sur le revenu après
cotisations, 5~\% au premier euro, 13~\% à 20~000~€, 26~\% à 30~000~€, puis monte doucement jusqu'à 50~\% à
250~000~€.

\figgrand{ircsg_x02_aplanissement}{Étape 2 : taux moyen et taux marginal}{fig:ex2}

\textbf{Ce qui change.} Les marches disparaissent (figure~\ref{fig:ex2}). Pour une personne seule à
50~000~€ brut, le taux marginal passe de 32 à 24~\%. L'effet redistributif est le même qu'aujourd'hui et la
perte sèche baisse d'un quart, à 13~Md€. Les transferts qui restent tiennent à l'harmonisation elle-même :
les pensions paient 10,9~Md€ de plus (1~430~€ pour un retraité seul à 17~000~€), les déciles 2 à 4 de 90 à
270~€ de plus, les déciles 5 à 7 de 280 à 430~€ de moins. Les couples à un seul revenu paient encore
1~540~€ de plus à 80~000~€ avec deux enfants. Les perdants de plus de 500~€ restent 13~\% des foyers.

\textbf{Le défaut : un barème proche d'aujourd'hui garde l'essentiel de son coût}, car ses taux marginaux de
25 à 35~\% portent sur la masse des revenus du travail, entre 30~000 et 120~000~€.

\subsection{Étape 3 : ce que donne la seule efficacité}

\textbf{La règle.} Parmi les mêmes barèmes, on retient celui qui minimise la perte sèche, à même recette,
sans aucune autre contrainte. C'est le meilleur des barèmes essayés, pas un optimum : il bute sur les bornes
de la grille, et sans taux supérieur imposé il tendrait vers un impôt proportionnel.

\textbf{Le résultat est un impôt presque proportionnel} (figure~\ref{fig:ex3}) : 14~\% au premier euro,
16~\% à 20~000~€, 18~\% à 120~000~€, et 50~\% n'est atteint qu'à 600~000~€ par adulte, la limite des
barèmes essayés. La perte sèche tombe à 5~Md€ (2 à 9 selon les élasticités), plus de trois fois moins
qu'aujourd'hui. Le mécanisme est simple : la perte croît avec le carré du taux marginal et les hauts revenus
réagissent le plus ; on abaisse donc leur taux marginal et on reprend la recette sur tout le monde, en
particulier sur les pensions, qui ne réagissent pas.

\figgrand{ircsg_x03_perte_seche}{Étape 3 : taux moyen et taux marginal}{fig:ex3}

\textbf{Le défaut : une forte redistribution à l'envers.} Les déciles 2 à 4 paient 610 à 1~310~€ de plus par
an (6 à 8~points de revenu), le dernier décile 5~380~€ de moins ; 57~\% des foyers perdent plus de 500~€. Un
salarié seul au SMIC paie 800~€ de plus. Et le gain va surtout aux personnes seules aisées : à 80~000~€ brut,
une personne seule gagne 7~120~€, un couple avec deux enfants 1~030~€, soit 7,6~points de revenu d'écart.

\textbf{À quel point les gens réagissent-ils vraiment ?} Tout repose sur les élasticités. À court terme, les heures
de la plupart des salariés bougent peu, en partie parce que les horaires sont fixés par l'employeur, ce qui
masque une réaction de long terme plus forte (Chetty, Friedman, Olsen et Pistaferri, 2011) ; une année
blanche d'impôt en Suisse n'a guère changé l'emploi (Martínez, Saez et Siegenthaler, 2021). Les réactions sont nettes chez les hauts revenus et les indépendants, qui changent
surtout la forme ou la déclaration de leurs revenus (Cabannes, Houdré et Landais, 2014 ; Sicsic, 2022),
chez les seconds apporteurs et chez les seniors. Entre pays, le lien entre taxation du travail et heures
disparaît quand on tient compte des dépenses sociales et du droit du travail (Gethin et Saez, 2026). Avec
les élasticités basses, le gain d'efficacité de l'étape 3 n'est que de 6~Md€ par an ; avec les hautes,
de 22~Md€.

\subsection{Étape 4 : notre proposition}

On part du barème de l'étape 3 et on corrige ses deux défauts, en acceptant de perdre une partie du gain
d'efficacité. Le barème est à chaque fois réoptimisé : c'est toujours celui qui minimise la perte sèche
sous les nouvelles règles.

\textbf{Première correction : une franchise de 8~000~€ par adulte sur les revenus d'activité.} Les
premiers 8~000~€ de salaire, de bénéfice ou de loyer de chaque adulte ne sont pas imposés ; les pensions et
les allocations chômage n'en bénéficient pas. Elle fait deux choses. Pour un salarié au SMIC, elle compense
la hausse du taux d'entrée : il paie ce qu'il paie aujourd'hui. Et elle rend le premier emploi, ou le
retour à l'emploi d'un conjoint, plus payant, puisque ses premiers euros ne sont pas imposés. Le
tableau~\ref{tab:franchise} montre le prix de ce choix : chaque millier d'euros de franchise oblige à relever
les taux et coûte environ 0,5~Md€ de perte sèche. Au-delà de 8~000~€, la hausse des déciles 2 à 4 ne baisse
plus : le taux d'entrée qu'il faut relever pour financer la franchise reprend ce qu'elle donne.

\begin{table}[!htbp]\centering\scriptsize\setlength{\tabcolsep}{3pt}
\caption{Le prix de la franchise sur les revenus d'activité (enfants : franchise de 6~000~€)}\label{tab:franchise}
\begin{tabular}{@{}lrrrrr@{}}
\toprule
Franchise par adulte & 0~€ & 4~000~€ & 6~000~€ & \textbf{8~000~€} & 10~000~€ \\
\midrule
Perte sèche (Md€ par an) & 5,2 & 7,7 & 8,3 & \textbf{9,3} & 10,1 \\
Déciles 2 à 4 (€ par foyer et par an) & +610 à +1~310 & +170 à +790 & +130 à +570 & \textbf{+210 à +450} & +290 à +390 \\
Foyers qui perdent plus de 500~€ & 57~\% & 43~\% & 37~\% & \textbf{33~\%} & 32~\% \\
Salarié seul au SMIC (€ par an) & +800 & +340 & +120 & \textbf{$-$10} & $-$190 \\
Retraité seul, 17~000~€ (€ par an) & +2~520 & +2~740 & +2~960 & \textbf{+3~250} & +3~540 \\
\bottomrule
\end{tabular}
\end{table}

\textbf{Seconde correction : la famille.} La première correction rapproche déjà familles et personnes
seules, parce qu'elle oblige à relever les taux du milieu du barème, dont profitaient surtout les personnes
seules aisées. Reste à choisir comment l'impôt tient compte des enfants. Des quatre traitements comparés à la
section suivante, la franchise de 6~000~€ par enfant, déjà utilisée aux étapes 1 à 3, est celle qui
rapproche le plus les familles des personnes seules : à 80~000~€ brut, la personne seule gagne 4~410~€, le
couple avec deux enfants 1~630~€, soit 3,5~points de revenu d'écart, contre 7,6 à l'étape 3 et 7,5 sans aucune
aide familiale dans l'impôt.

\textbf{Le barème proposé} prélève, par adulte et sur le revenu après cotisations et franchise, 17~\% au
premier euro, 22~\% à 20~000~€, 24~\% à 45~000~€, puis monte en ligne droite jusqu'à 50~\% à 600~000~€
(figure~\ref{fig:prop}). La perte sèche est de 9~Md€ (4 à 16 selon les élasticités), deux fois moins
qu'aujourd'hui. Ce gain doit être lu avec deux réserves. Il vient pour l'essentiel (plus de 90~\%) du
dixième le plus aisé, dont l'élasticité mesure surtout des changements de forme ou de calendrier des revenus,
et pour partie des passages vers l'impôt sur les sociétés ou le prélèvement forfaitaire : une partie du gain
est un déplacement entre impôts plutôt qu'un surcroît d'activité. Et le calcul ne retient que l'IR et la CSG
dans le coin fiscal ; avec les cotisations non contributives et la prime d'activité, les pertes seraient plus
élevées à tous les niveaux et l'ordre des barèmes pourrait changer. Un tiers des foyers perdent plus de 500~€, contre 57~\% à l'étape 3.

\figgrand{ircsg_x04_proposition}{La proposition : taux moyen et taux marginal}{fig:prop}

\textbf{Qui finance.} Par type de revenu, les salaires paient 13,3~Md€ de moins, les revenus des
indépendants et les loyers 2,6~Md€ de moins ; les pensions paient 13,0~Md€ de plus, les allocations chômage
3,0~Md€. La fin des niches élargit l'assiette, mais c'est l'imposition des retraités au même barème que les actifs, quel
que soit leur revenu, qui finance l'essentiel de la baisse sur le travail. Un retraité seul à 17~000~€ de
pension, qui ne paie rien aujourd'hui, paierait 3~250~€ ; à 30~000~€, 2~830~€ de plus. Comme la franchise ne
porte que sur les revenus d'activité, un retraité modeste paie même nettement plus qu'un salarié de même
revenu brut : à 17~000~€, 3~250~€ contre 1~270~€ ; à 30~000~€, 6~140~€ contre 3~600~€. C'est un choix : la
franchise récompense l'activité, pas le revenu ; dans le modèle, elle couvre aussi les loyers, qu'il serait
plus cohérent d'en exclure. Une limite s'impose en bas : à 10~000~€ de pension, l'impôt serait de 1~800~€, et
les titulaires du minimum vieillesse (12~523~€ par an) tomberaient sous ce minimum. Il faut exonérer les
pensions jusqu'à ce niveau, ce que le modèle ne fait pas.

\textbf{Les très hauts revenus.} Leurs revenus du capital restent au prélèvement forfaitaire, majoration
de 3 puis 4~points comprise (note PFU) ; seule change l'imposition de leurs revenus du barème. Le résultat
dépend de la part de leurs revenus qui relève du prélèvement forfaitaire, que la DGFiP ne publie pas par
tranche. Dans la convention centrale du modèle, les 3~800 foyers déclarant plus de 3~M€ paieraient 3 à
6~points de revenu de plus ; mais cette convention prête à leurs revenus du barème un IR actuel d'environ
23~\%, bien en dessous de leur taux marginal de 49~\%, ce qui gonfle la hausse. Si l'on suppose plutôt que le
capital se concentre en haut, ce que suggère la concentration des dividendes, les foyers de 250\,000~€ à
9~M€ paient 2 à 6,5~points de moins, le taux marginal supérieur passant de 55 à 50~\%, et les 700 foyers
au-delà autant qu'aujourd'hui ; le barème doit alors être un peu plus haut (25,4~\% au lieu de 24,2~\% à
45\,000~€), le reste bougeant peu (perte sèche de 9,0~Md€, 34~\% de perdants). Caler cette part sur les
données d'impôt par tranche est la prochaine étape du modèle.

\subsection{La famille : quatre façons de faire}\label{sec:famille}

Aujourd'hui, le quotient familial des enfants réduit l'IR d'environ 14~Md€ : 13,9~Md€ dans notre modèle
pour les revenus 2024, dont 9,7 pour les couples et 4,2 pour les parents isolés. L'Insee trouvait 13,3~Md€
pour les familles avec enfants en 2017, sur 16,6~Md€ pour l'ensemble des demi-parts familiales (André et
Sireyjol, 2021). Le tableau~\ref{tab:famille} compare quatre traitements, avec la même franchise
d'activité de 8~000~€ et un barème à chaque fois réoptimisé.

\begin{table}[!htbp]\centering\scriptsize\setlength{\tabcolsep}{3.5pt}
\caption{Quatre traitements de la famille dans l'impôt fusionné (franchise d'activité de 8~000~€)}\label{tab:famille}
\begin{tabular}{@{}lrrrr@{}}
\toprule
& Rien & \textbf{Franchise} & Chèque & Quotient \\
\midrule
Aide aux familles dans l'impôt (Md€ ; aujourd'hui 13,9) & 0 & \textbf{20,6} & 13,9 & 3,6 \\
Perte sèche (Md€ par an) & 7,8 & \textbf{9,3} & 8,2 & 8,1 \\
Foyers qui perdent plus de 500~€ & 33~\% & \textbf{33~\%} & 36~\% & 31~\% \\
\midrule
\multicolumn{5}{@{}l}{\emph{Variation moyenne par foyer et par an}} \\
Personnes seules sans enfant & $-$200~€ & \textbf{+130~€} & +50~€ & $-$90~€ \\
Parents isolés & +1~030~€ & \textbf{+340~€} & +510~€ & +900~€ \\
Couples sans enfant & $-$500~€ & \textbf{+360~€} & +20~€ & $-$450~€ \\
Couples avec deux enfants & +530~€ & \textbf{$-$970~€} & $-$500~€ & +190~€ \\
\midrule
\multicolumn{5}{@{}l}{\emph{Foyers types, 80~000~€ brut}} \\
Personne seule & $-$5~910~€ & \textbf{$-$4~410~€} & $-$5~020~€ & $-$4~970~€ \\
Couple, deux enfants, 60/40 & +70~€ & \textbf{$-$1~630~€} & $-$520~€ & $-$240~€ \\
Couple, deux enfants, un seul revenu & +2~560~€ & \textbf{+1~080~€} & +1~900~€ & +1~300~€ \\
\bottomrule
\end{tabular}
\par\vspace{3pt}\parbox{\linewidth}{\raggedright\scriptsize\color{encre}
Rien : aucune aide familiale dans l'impôt. Franchise : 6~000~€ par enfant déduits du revenu du parent qui
gagne le plus. Chèque : le coût actuel du quotient familial des enfants reversé en un montant égal par
enfant (772~€ par an), versé même sans impôt. Quotient familial : quotient conjugal et familial appliqués à
l'impôt fusionné, plafond de 1~791~€ par demi-part comme aujourd'hui. Aide aux familles : ce que les foyers
avec enfants paieraient de plus s'ils n'en avaient pas.}
\end{table}

\textbf{Rien dans l'impôt.} C'est alors aux prestations de prendre le relais. Les familles perdent : 530~€ en moyenne
pour un couple avec deux enfants, 1~030~€ pour un parent isolé, et l'écart avec les personnes seules
atteint 7,5~points de revenu à 80~000~€. C'est la variante la moins coûteuse en efficacité, parce qu'elle
n'a rien à financer.

\textbf{Un chèque égal par enfant.} Reverser le coût actuel du quotient familial donne 772~€ par enfant et
par an. Il faut pour cela que l'impôt rapporte 13,9~Md€ de plus avant chèques, 234~Md€ au lieu de 220. Les
personnes seules sans enfant paient-elles plus ? Par rapport à aujourd'hui, à peine : 50~€ par foyer en
moyenne, puisqu'elles finançaient déjà le quotient. Par rapport à une fusion sans aide familiale, oui :
250~€ de plus. Le chèque donne autant à chaque enfant, quel que soit le revenu des parents : les familles
aisées y gagnent moins qu'avec le quotient d'aujourd'hui, d'où un gain moyen limité pour les couples avec
enfants (500~€).

\textbf{Le quotient familial sur l'impôt fusionné.} Le quotient ne vaut que par la progressivité : il
divise le revenu pour le faire passer sous des taux plus bas. Or le barème fusionné est presque
proportionnel entre 20~000 et 100~000~€ : appliqué à lui, le quotient ne vaut plus que 3,6~Md€ pour les
enfants, et les couples avec deux enfants paient 190~€ de plus en moyenne. Le quotient conjugal, lui, profite
surtout aux couples sans enfant (450~€ de moins).

\textbf{La franchise par enfant, retenue,} est la plus favorable aux familles : elle vaut de 1~300 à
1~550~€ par enfant et par an selon le revenu du parent, plus que le quotient d'aujourd'hui pour les familles
modestes et moyennes, et davantage aussi au-delà d'environ 170\,000~€ de revenu, car sa valeur suit le taux
marginal alors que le quotient est plafonné à 1~791~€. Les familles de trois enfants, que le quotient
favorise aujourd'hui, ne sont pas modélisées. Au total, 20,6~Md€ au lieu de 13,9 (16,2 en comptant les
réductions pour frais de garde et de scolarité, supprimées avec les niches), financés par des taux un peu plus élevés pour tous, d'où une perte sèche
plus forte de 1,1 à 1,5~Md€ que dans les autres variantes. C'est le prix de la seconde correction.

\subsection{Qui perd, qui gagne, à chaque étape}

\begin{table}[!htbp]\centering\scriptsize\setlength{\tabcolsep}{3pt}
\caption{Les étapes en chiffres}\label{tab:synth}
\begin{tabular}{@{}lrrrrr@{}}
\toprule
& Aujourd'hui & Étape 1 & Étape 2 & Étape 3 & \textbf{Proposition} \\
\midrule
Taux marginal, seul à 50~000~€ brut & 31,9~\% & 29,6~\% & 24,4~\% & 14,8~\% & \textbf{21,1~\%} \\
Perte sèche (Md€ par an) & 18 & 18 & 13 & 5 & \textbf{9} \\
Effet redistributif (aujourd'hui = 100) & 100 & 126 & 97 & 21 & \textbf{75} \\
\midrule
Foyers qui perdent plus de 500~€ & --- & 14~\% & 13~\% & 57~\% & \textbf{33~\%} \\
Foyers qui perdent plus de 1~000~€ & --- & 11~\% & 8~\% & 32~\% & \textbf{11~\%} \\
Foyers qui gagnent plus de 500~€ & --- & 38~\% & 21~\% & 17~\% & \textbf{16~\%} \\
\midrule
Déciles 2 à 4 (€ par foyer) & --- & $-$80 à $-$150 & +90 à +270 & +610 à +1~310 & \textbf{+210 à +450} \\
Dernier décile (€ par foyer) & --- & +3~040 & +760 & $-$5~380 & \textbf{$-$1~440} \\
Pensions (Md€) & --- & +8,8 & +10,9 & +16,7 & \textbf{+13,0} \\
Retraité seul, 17~000~€ & --- & +1~110~€ & +1~430~€ & +2~520~€ & \textbf{+3~250~€} \\
Salarié seul au SMIC & --- & $-$650~€ & $-$370~€ & +800~€ & \textbf{$-$10~€} \\
Personne seule, 80~000~€ & --- & $-$1~910~€ & $-$3~390~€ & $-$7~120~€ & \textbf{$-$4~410~€} \\
Couple, 2 enfants, 80~000~€, 60/40 & --- & $-$4~130~€ & $-$2~880~€ & $-$1~030~€ & \textbf{$-$1~630~€} \\
Couple, 2 enfants, 80~000~€, un revenu & --- & +2~590~€ & +1~540~€ & $-$690~€ & \textbf{+1~080~€} \\
\bottomrule
\end{tabular}
\par\vspace{3pt}\parbox{\linewidth}{\raggedright\scriptsize\color{encre} Foyers types sans niche (OpenFisca) ;
déciles, pensions et perdants : foyers réels, niches comprises dans la charge actuelle. Les foyers types
paient souvent plus aujourd'hui qu'un foyer réel de même revenu, d'où des gains plus élevés que dans les
moyennes par décile.}
\end{table}

\subsection{Comment passer d'un système à l'autre}

La collecte ne change presque pas. Les employeurs, les caisses de retraite et France Travail prélèvent déjà à
la source l'IR et la CSG sur les mêmes revenus ; ils appliqueraient un seul taux au lieu de deux. Le taux
individualisé dans les couples est même devenu la règle par défaut du prélèvement à la source depuis le
1\textsuperscript{er}~septembre~2025 (loi de finances pour 2024, art.~19) : l'impôt individuel n'y ajoute
qu'un calcul de franchise par adulte.

Le vrai sujet est le rythme : un retraité à 17~000~€ de pension absorberait difficilement 3~250~€ en une
année.
Une convergence linéaire sur quatre ans, où chaque paramètre (abattements, taux de CSG, barème) parcourt
chaque année un quart du chemin, limite la hausse annuelle à environ 800~€ pour ce retraité, sans changer
la cible. Comme les gagnants et les perdants avancent au même rythme, la recette reste à peu près
constante chaque année.

\textbf{Le risque constitutionnel.} Le Conseil constitutionnel a censuré en 2000 une réduction de CSG qui
ignorait les autres revenus du foyer et les charges de famille (décision 2000-437 DC), et en 2012 une
contribution assise sur chaque personne sans tenir compte du foyer (2012-662 DC). Un impôt individuel avec
une franchise réservée aux revenus d'activité réunit ces deux traits. La proposition s'en distingue parce
qu'elle remplace l'impôt de foyer au lieu de s'y ajouter et qu'elle tient compte des enfants, mais
l'argument reste à construire ; la variante « quotient », qui rétablit l'imposition commune des couples, est
la solution de repli (perte sèche de 8,1~Md€).


\section{Immobilier : moins taxer les ventes, mieux taxer le sol}\label{sec:foncier}
% Partie foncier (source : ../foncier/note-foncier.tex, corrigée selon ../revue/verif_foncier.md)
La France taxe beaucoup l'immobilier, et elle le taxe au mauvais endroit. Un ménage qui achète un logement
ancien paie 5,8 à 6,3~\% du prix en droits de mutation, comptant, au moment où son apport est le plus
contraint. La taxe foncière, elle, repose sur des valeurs locatives de 1970, qui font payer
proportionnellement plus les logements des communes pauvres que ceux des communes chères. Un impôt sur les
ventes décourage les ventes : il retient dans leur logement des ménages qui voudraient en changer et pèse
d'abord sur ceux qui achètent. Un impôt récurrent sur la valeur du sol ne décourage ni le travail ni la
construction, puisque le terrain existe de toute façon ; son défaut est politique : il se voit, chaque
année, sur un avis d'imposition.

\subsection{État des lieux : 116 milliards, dont un cinquième sur les ventes}

Les prélèvements assis sur l'immobilier atteignent environ 116~Md€ en 2024, 4~\% du PIB
(figure~\ref{fig:inv}). La détention en fournit 69~Md€ : taxe foncière sur les propriétés bâties
(42~Md€, dont environ 31,6~Md€ sur les logements), taxe d'enlèvement des ordures ménagères (9~Md€),
cotisation foncière des entreprises (8~Md€), impôt sur la fortune immobilière (2,7~Md€) et taxes
d'habitation restantes. Les transactions en fournissent 25~Md€ : droits de mutation (14,1~Md€ en 2024,
16,7 en 2025, dont 13,6 sur les logements), TVA sur le neuf et les terrains (7,1~Md€), plus-values
immobilières (2,8~Md€), contribution de sécurité immobilière (0,7~Md€). S'y ajoutent les successions et
donations portant sur l'immobilier (environ 12,5~Md€, estimation), l'impôt sur les revenus locatifs
(8,4~Md€) et la taxe d'aménagement (1,5~Md€). Selon l'OCDE, la France prélevait en 2023 1,90~\% du PIB en
impôts périodiques sur l'immobilier, le double de la moyenne (0,95~\%), et 0,72~\% sur les transactions,
contre 0,37~\%. Elle taxe donc lourdement les deux ; l'argument n'est pas que la détention serait peu taxée,
mais qu'un euro de droits de mutation coûte davantage à l'économie qu'un euro de taxe foncière, surtout
s'il porte sur le sol.

\figbloc{fon_03_inventaire}{Les prélèvements assis sur l'immobilier en 2024}{fig:inv}

\textbf{Les droits de mutation} d'un logement ancien cumulent la part départementale (4,50~\%), la taxe
communale (1,20~\%) et les frais d'assiette de l'État, soit 5,81~\% du prix. La loi de finances pour 2025 a
permis aux départements de passer à 5~\% jusqu'au 31~mars~2028 ; 88 l'ont fait (6,32~\% au total), sauf pour
les primo-accédants. \textbf{La taxe foncière} repose sur la moitié de la valeur locative de 1970,
revalorisée par des coefficients nationaux ; la révision des valeurs des logements a été reportée à 2031.
Rapportée à la valeur de marché, elle pèse 0,52~\% dans le dixième des logements situés dans les communes
les moins chères et 0,18~\% dans le dixième le plus cher (figure~\ref{fig:tfv}) ; à Paris 0,12~\%, malgré la
hausse du taux de 13,5 à 20,5~\% en 2023. Elle est régressive rapportée au patrimoine et progressive rapportée
au revenu (André et Meslin, 2023).

\subsection{Ce que dit la recherche}

\textbf{Les droits de mutation réduisent les ventes et la mobilité.} Les ventes réagissent d'abord au
calendrier : le mois qui précède une hausse, elles bondissent de 26~\% en France en 2014 (Bérard et Trannoy,
2018). Au-delà, la plupart des études étrangères trouvent qu'un point de droits supprime environ 8~\% des
ventes, dans une fourchette de 0 à 18~\% et sur des horizons de quelques mois à quatre ans (Dachis, Duranton
et Turner, 2012 ; Fritzsche et Vandrei, 2019 ; Best et Kleven, 2018), et que la mobilité des propriétaires
baisse de 8 à 14~\% par point. En Finlande, ce sont surtout les déménagements vers plus grand qui
disparaissent (Eerola et al., 2021) : la taxe frappe le ménage qui veut une pièce de plus, souvent parce que
la famille s'agrandit, ce que l'étude n'observe pas directement. En France, les preuves sont partagées :
après la hausse de 0,7~point de 2014, les ventes baissent de 6~\% sur trois mois puis plus rien selon Bérard et
Trannoy, de 10 à 12~\% durablement selon Bachelet et Poulhès (2017), qui trouvent que l'acheteur a supporté
toute la hausse ; en Allemagne, un point de droits réduit les prix d'environ 3~\%, et ce sont les vendeurs qui
paient (Dolls, Fuest, Krolage et Neumeier, 2025). Le seul modèle français (Bontemps, Cherbonnier et Magnac,
2023) trouve un gain modeste à remplacer les droits par un impôt foncier, qui s'inverse si la baisse est
financée par l'impôt sur le revenu ; la baisse y profite surtout aux jeunes ménages aisés, dont l'accès à la
propriété progresse, alors que celui des jeunes ménages modestes recule.

\textbf{La taxe foncière : un bon impôt mal assis.} Une hausse permanente se capitalise dans les prix : le
propriétaire présent à l'annonce perd la valeur actuelle des hausses futures, l'acheteur suivant paie moins
cher (Koster et Pinchbeck, 2022). Pour un impôt sur le seul sol, l'ampleur de cette capitalisation reste
discutée (Høj et al., 2018 ; Nielsson et al., 2024). Les évaluations de masse sont souvent régressives,
parce qu'elles suivent mal les quartiers chers (Berry, 2021 ; Avenancio-León et Howard, 2022) ; Pittsburgh a
abandonné son double taux en 2001 après une réévaluation brutale, Philadelphie a réussi la sienne en 2014
avec un abattement pour les occupants de longue date. Un plafond de hausse attaché au propriétaire crée un
verrou (Proposition~13 en Californie ; Ferreira, 2010). Le double taux, plus lourd sur le terrain, a un effet
positif mais modeste sur la construction : en Pennsylvanie, un rapport de quatre pour un a accru d'environ
3~points la croissance décennale du nombre de pièces (Banzhaf et Lavery, 2010). Le Bade-Wurtemberg taxe
depuis 2025 le seul sol, par des valeurs de zone ; la plus haute juridiction fiscale allemande l'a validé en
mai 2026, des recours constitutionnels restant pendants. À Karlsruhe, à recette constante, les maisons
paient 76~\% de plus et les copropriétés 20~\% de moins.

\subsection{La bascule et ses arbitrages}

\begin{encadre}{Comment la bascule est chiffrée}
\textbf{Champ.} Les 30,5~millions de logements des particuliers (parc du recensement moins les HLM) de France
métropolitaine hors Alsace-Moselle. Valeur : prix médian au m² des ventes 2023-2024 de la commune multiplié
par la surface médiane des logements vendus, soit 6~760~Md€, 17 à 19~\% sous les comptes de patrimoine de
l'Insee : les recettes sont justes, mais les taux exprimés en part de la valeur sont surestimés d'autant
(0,19~\% sur nos valeurs correspond à environ 0,15~\% sur celles de l'Insee).

\textbf{Part du terrain.} Maisons : prix des terrains à bâtir multiplié par la surface de la parcelle ;
appartements : valeur moins coût de reconstruction, calé pour retrouver la part du terrain des comptes de
patrimoine (45~\%). Elle va de 16~\% dans le dixième le moins cher à 67~\% dans le plus cher.

\textbf{Charges actuelles.} Taxe foncière : rapport taxe/valeur de chaque commune (Insee, 2017), vieilli selon
la croissance de sa taxe foncière de 2017 à 2024 (REI), calé sur 23,6~Md€ ; droits : 5,81~\% des ventes de
logements anciens, 11,4~Md€ par an en 2023-2024, creux du cycle.

\textbf{Non modélisé.} Capitalisation, âge et revenu des propriétaires, locaux professionnels, résidences
secondaires.
\end{encadre}

\textbf{La référence.} Les droits de mutation sur les logements sont divisés par deux, de 5,81 à 2,9~\%
(5,7~Md€ si les ventes ne réagissent pas, 4,6~Md€ si elles augmentent de 6~\% par point). Pour éviter que les
ventes s'arrêtent dans l'attente, la baisse s'applique dès son annonce, comme en Finlande en 2023. Chaque
commune garde sa recette de taxe foncière et son pouvoir de taux, mais l'assiette devient la valeur de
marché, avec un taux sur le terrain égal à quatre fois le taux sur le bâti. Les droits perdus sont
remplacés par un taux national de 0,19~\% sur la valeur des terrains, calé sans reprise des ventes. L'impôt
qui porte sur les constructions passe d'environ 14 à 7~Md€ : c'est le principal effet économique.

\figbloc{fon_01_tf_valeur}{La taxe foncière rapportée à la valeur des logements, aujourd'hui et après la bascule}{fig:tfv}

\textbf{Résultats.} La taxe foncière augmente de 3~\% dans le dixième des logements le moins cher et de 69~\%
dans le plus cher. En comptant la baisse des droits, rapportée à l'ensemble des logements de la commune, la
charge varie de $-$75 à +133~€ par logement et par an selon le décile (figure~\ref{fig:var}). Rapportée à la
valeur, la taxe va de 0,54~\% dans les communes les moins chères à 0,30~\% dans les plus chères : l'écart se
resserre sans être renversé. Les hausses se concentrent à Paris. Un logement sur quatre voit sa taxe
foncière augmenter de plus de 300~€ ; un sur vingt la voit plus que doubler, pour 70~\% à Paris et 14~\% dans
les Hauts-de-Seine. À Paris, la hausse moyenne est de 122~\% ; hors Paris, le dixième le plus cher monte de
51~\%. Le partage entre maisons et appartements n'est pas robuste : avec notre mesure du terrain, les
appartements paient davantage (+240~€ contre +150~€) ; si l'on mesure le terrain des maisons comme celui des
appartements, c'est l'inverse (+130~€ contre +230~€). Il faudra simuler l'assiette légale, valeur de zone
multipliée par la parcelle et répartie entre copropriétaires, avant de conclure.

\figbloc{fon_02_variation}{Variation de la charge par logement et par an, selon le prix au m² de la commune}{fig:var}

\textbf{La perte sèche.} Chaque vente découragée est un échange manqué, qui valait au plus les droits
payés : le gain de la baisse est le triangle sur les ventes, $\sigma V (\tau_0^2-\tau_1^2)/2$, avec 196~Md€ de
ventes par an et $\sigma$ la baisse des ventes par point. À 6~\% par point, il vaut 1,5~Md€ par an (0,8 à
3,7~Md€ pour 3 à 15~\%). Le terrain ne réagit pas à l'impôt ; le bâti, lui, est moins taxé : 0,4~Md€ de plus
si le capital logement réagit proportionnellement au coût d'usage. Le gain central est donc d'environ
2~Md€ par an, 0,07~\% du PIB. Il pourrait doubler si l'on compte le coût des logements mal adaptés que la
mobilité réduite laisse en place (de l'ordre de 3,5~Md€ selon le modèle finlandais de Määttänen et Terviö,
2022, non calibré sur la France) ; il pourrait aussi être plus faible, parce que le logement occupé par son
propriétaire est déjà avantagé (le loyer qu'il se verse à lui-même n'est pas imposé), si bien qu'alléger
l'impôt sur le bâti l'avantage davantage.

\textbf{Les arbitrages} (tableau~\ref{tab:fon_arb}).

\begin{table}[!htbp]\centering\scriptsize
\caption{Les arbitrages de la bascule}\label{tab:fon_arb}
\begin{tabularx}{\linewidth}{@{}>{\raggedright\arraybackslash}p{3.3cm}>{\raggedright\arraybackslash}X>{\raggedright\arraybackslash}p{4.3cm}@{}}
\toprule
Choix & Ce qu'il change & Qui paie \\
\midrule
Baisse des droits : 50 ou 30~\% & 50~\% : 5,7~Md€ à reporter, perte sèche $-$1,5~Md€ ; 30~\% : 3,4~Md€ (2,3 si les ventes réagissent), $-$1,0~Md€ & taxe foncière +24 ou +14~\% en moyenne \\
Où reporter : taxe actuelle relevée & Ne change rien aux valeurs de 1970 & +24~\% partout ; 0,65~\% de la valeur dans les communes pauvres, 0,22~\% dans les riches \\
Où reporter : deux taux nationaux & Supprime d'un coup les écarts de 1970 & dernier décile +215~\% ; un logement sur quatre à plus de +50~\% \\
Où reporter : deux taux par commune et taux national sur le terrain (référence) & Garde la recette et le pouvoir de taux des communes & +3 à +69~\% selon le décile ; Paris +122~\% \\
Rapport terrain/bâti : 2, 4 ou sol seul & Plus le rapport est élevé, plus le bâti est allégé ($-$0,2, $-$0,4, $-$0,5~Md€ de perte sèche) & avec deux taux nationaux, dernier décile +184, +215, +266~\% ; maisons +10, +4, $-$6~\% \\
Taux national calé sans ou avec reprise des ventes & 0,19 ou 0,15~\% ; écart de 1,1~Md€ & à 0,19~\%, les propriétaires portent le risque ; à 0,15~\%, le budget \\
Compensation des collectivités & Taux national reversé aux départements ; la taxe communale de 1,2~Md€ reste à compenser & Paris paie 213~M€ de plus qu'elle ne perd, les Hauts-de-Seine 106~M€ (part départementale seule) \\
Report de paiement des propriétaires âgés & Facultatif ou proposé par défaut, avec intérêts & recette différée \\
\bottomrule
\end{tabularx}
\end{table}

\figbloc{fon_04_primo}{Ce que gagne un ménage qui achète le logement médian de sa commune}{fig:primo}

\textbf{Qui gagne.} Le gain d'un jeune acheteur dépend de qui supporte l'impôt (figure~\ref{fig:primo}). Si
les prix ne bougent pas, comme après la hausse de 2014, l'acheteur du logement médian paie entre 2~300~€
(appartement à Saint-Étienne) et 22~700~€ (maison à Arcachon) de droits en moins : 8~100~€ à Lyon, 13~000~€
pour une maison à Nantes, 12~100~€ pour un appartement à Paris. La hausse de taxe foncière qu'il paie ensuite,
actualisée à 3~\%, ne rattrape cette économie qu'au bout de 24 à 39~ans à Paris, Lyon, Bordeaux et Limoges, et
jamais à Nantes, Arcachon, Ussel ou Guéret. Le perdant est le propriétaire installé qui ne vend pas : jusqu'à
660~€ par an pour l'appartement médian parisien, 22~000~€ en valeur actuelle. Si au contraire les deux
mouvements passent dans les prix, le prix moyen bouge peu, et la bascule transfère de la valeur entre
logements (des plus fonciers et des moins souvent vendus vers les autres) plutôt qu'entre générations :
l'appartement médian parisien perdrait environ 8~000~€, 2~\% de sa valeur. Dans les deux cas, l'argument le
plus solide pour les jeunes est la liquidité : si l'apport représente 10 à 20~\% du prix, la baisse de
2,9~points en représente un septième à un tiers. Pour les familles, la baisse divise par deux le coût fiscal
d'un agrandissement ; aucune étude ne la relie directement à la natalité.

\subsection{Comment passer d'un système à l'autre}

La France n'a pas de valeurs du sol par zone, comme les \emph{Bodenrichtwerte} allemandes ; les ventes
publiées (DVF, DV3F) permettent de les établir, et le Bade-Wurtemberg a mis cinq ans à le faire. La baisse
des droits ne peut pas attendre : annoncée pour plus tard, elle gèlerait les ventes. Elle s'applique donc
dès son vote, financée le temps d'établir les valeurs par une hausse uniforme et temporaire de la taxe
foncière actuelle (20 à 24~\%), remplacée ensuite par la nouvelle assiette, qui se substitue à la révision
reportée à 2031. Le taux national sur le terrain s'applique en totalité dès la première année ; seule la
nouvelle répartition de la taxe communale est étalée sur dix ans, hausses et baisses au même rythme, pour
que la recette de chaque commune reste constante. Les valeurs seraient révisées tous les trois ans et
lissées, contestables par expertise au-delà d'un écart de 30~\% ; une révision à chaque vente, comme celle
adoptée en commission le 9~octobre~2026 puis tombée avec la première partie du projet de loi de finances
(amendement I-CF1872), recréerait le verrou de la Proposition~13. Le Conseil constitutionnel n'exige pas de
compensation à l'euro près d'une recette locale supprimée (décision 2019-796 DC), mais les départements
échangeraient la moitié de leur principal impôt à taux voté contre une dotation, et la taxe d'ordures
ménagères, assise sur la même base, devrait en être découplée. Le seuil d'exonération des propriétaires de
plus de 75~ans, fondé sur le revenu fiscal de référence, devra être redéfini si l'IR et la CSG fusionnent.


\section{Impôt sur les sociétés : ne pas toucher au taux, mieux déduire l'investissement}\label{sec:is}
% Partie IS (source : ../is/note-is.tex, corrigée selon ../revue/verif_is_pfu.md ; chiffres : ../is/sorties/simulation_is.txt, euros de 2027)
L'impôt sur les sociétés (IS) n'appelle pas de grande réforme. Ce que lui reprochent les comparaisons
internationales, c'est son taux affiché, le plus élevé de l'OCDE pour les plus grands groupes en 2025 et 2026
avec la contribution exceptionnelle, sa complexité et son instabilité, pas le traitement de
l'investissement, qui est dans la moyenne. La référence garde donc le taux normal à 25~\% ; la question est
de savoir s'il faut déduire l'investissement plus vite, et comment le payer.

\subsection{Ce que les comparaisons reprochent à l'IS français}

La Tax Foundation classe la France dernière des 38 pays de l'OCDE en 2025 et en 2026, au total comme pour
l'IS, à cause du taux de 36,1~\% des plus grands groupes (24,2~\% en moyenne), de quatre taux différents
(2,3 en moyenne), d'une surtaxe, d'une boîte à brevets et de la troisième subvention implicite à la
recherche la plus généreuse. Cet indice pénalise par construction toute incitation ciblée. Il ne pénalise
pas la déduction de l'investissement : la France est 15\up{e} sur 38 en 2025 et 16\up{e} en 2026 pour la
récupération des coûts, et les déductions d'une machine y valent 88~\% de son prix en valeur actuelle, contre
86~\% en moyenne. Les taux effectifs de l'OCDE vont dans le même sens. Le taux moyen effectif, qui compte pour
choisir où implanter un projet rentable, est de 23,6~\% contre 22,2~\% (13\up{e} sur 38) ; le taux marginal
effectif, qui compte pour décider d'un investissement de plus, est de 15,0~\% contre 16,5~\%, et de 9,8~\%
contre 10,4~\% pour les équipements (chiffres de 2025, qui semblent ignorer la contribution exceptionnelle).
Le Conseil des prélèvements obligatoires (2023, 2025) et le Sénat (2026) font de l'instabilité le premier
grief : la contribution exceptionnelle, présentée comme temporaire, a été prorogée une fois, et son assiette
est la moyenne de l'IS de deux exercices, dont l'un était clos quand elle a été votée.

\subsection{Déduire l'investissement plutôt que baisser le taux}

Une entreprise achète une machine de 100~€ qui dure sept ans. Elle en déduit aujourd'hui le coût par
fractions ; actualisées à 5~\%, ces déductions valent 85 à 90~€ selon que l'amortissement est linéaire ou
dégressif. Si elle peut déduire les 100~€ l'année de l'achat, le total ne change pas, seulement sa date ; mais
l'État rend alors 25~€ à l'achat et prélève ensuite 25~\% de ce que la machine rapporte. Il devient un
associé, à condition que les pertes soient elles aussi déduites sans limite, ce que le plafonnement de leur
imputation empêche en partie. Un investissement qui rapporte tout juste le rendement normal du capital ne
paie plus d'impôt ; restent imposés les profits qui dépassent ce rendement et ceux du capital déjà installé.
Une baisse du taux allège au contraire tous les profits, y compris ceux des machines achetées il y a dix ans et
ceux des rentes, qui ne dépendent pas de la décision d'investir. Sur la réforme américaine de 2017, un point
de capital supplémentaire au bout de dix ans coûte environ 10~\% de la recette de l'IS par an avec une baisse
du taux et 2,5~\% avec la déduction immédiate (Chodorow-Reich, Smith, Zidar et Zwick, 2026). L'effet sur
l'investissement est établi : +10 et +17~\% d'investissement éligible pour les déductions partielles
américaines de 2001-2004 et 2008-2010, deux fois plus dans les petites entreprises et seulement chez celles
qui payaient l'impôt (Zwick et Mahon, 2017) ; +11~\% pour les entreprises moyennes britanniques (Maffini, Xing
et Devereux, 2019). Ces mesures étaient temporaires et ont en partie avancé des achats ; d'où l'importance
d'une mesure permanente. Il n'existe aucune évaluation causale du suramortissement français de 2015.

Le coût d'usage du capital, ce que doit rapporter une machine pour couvrir son financement, son usure et
l'impôt, s'écrit $c = (r+\delta)(1-\tau z)/(1-\tau)$, où $z$ est la valeur actuelle des déductions pour 1~€.
Avec $z$ entre 0,85 et 0,90, l'impôt le majore de 3,5 à 5~\% ; avec la déduction immédiate, plus du tout. Si
le stock d'équipements des sociétés imposables (environ les trois quarts des 541~Md€ des sociétés non
financières) réagit à peu près proportionnellement au coût d'usage (élasticité de 0,5 à 1,6), le PIB gagnerait
0,04 à 0,2~\% à long terme, 0,1~\% au centre. Ce gain de production n'est pas un gain de bien-être du même
montant : il rémunère d'abord le capital supplémentaire qui le produit. La perte sèche évitée, le triangle
sur ce capital, n'est que de 0,05 à 0,16~Md€ par an, parce que le coin supprimé est petit ; l'essentiel de
l'intérêt de la mesure tient à ce qu'elle coûte peu.

\subsection{Combien cela coûte, et comment le payer}

Les sociétés non financières ont investi 84~Md€ en 2025 en machines, équipements et matériels informatiques ;
on suppose que les trois quarts sont faits par des sociétés imposables, hypothèse non mesurée à laquelle le
coût est proportionnel. La première année (2027), une déduction immédiate totale coûterait 14 à 15,5~Md€ ;
le coût baisse ensuite, parce que les machines achetées après la réforme n'ont plus d'amortissement à
déduire. À partir de la huitième année ne reste qu'un coût permanent de 1,4 à 2,0~Md€ par an, qui vient
seulement de ce que l'investissement croît. En valeur actuelle, une année d'investissement coûte 1,8 à
2,6~Md€ au taux d'actualisation des entreprises et 1,2 à 1,7~Md€ à celui de l'État, soit 0,5 à 0,8 point de
taux normal (un point vaut 3,2~Md€ en 2027 une fois le taux réduit supprimé).

\figbloc{is_01_cout_deduction}{Coût annuel de la déduction immédiate des équipements et ressources pour la financer}{fig:cout}

\begin{table}[!htbp]\centering\scriptsize
\caption{Les arbitrages de l'IS (Md€ de 2027 par an)}\label{tab:is_arb}
\begin{tabularx}{\linewidth}{@{}>{\raggedright\arraybackslash}p{3.4cm}>{\raggedright\arraybackslash}X>{\raggedright\arraybackslash}p{4.2cm}@{}}
\toprule
Choix & Ce qu'il change & Qui paie \\
\midrule
Déduction immédiate : 50~\% ou 100~\% & 50~\% : 7 à 7,8 la 1\up{re} année, gain de PIB divisé par deux ; 100~\% : 14 à 15,5, puis 1,4 à 2 & --- \\
Calendrier : 100~\% dès 2027, ou 50~\% jusqu'en 2032 puis 100~\% & dette maximale (intérêts à 3,25~\%) : 17 à 45 ou 2 à 32 ; remboursée en 10 à 28 ou 4 à 25~ans & --- \\
Financement par l'assiette & Fin du taux réduit (2,2 à 3,1), CIR apprécié au niveau du groupe (0,6), fin des exonérations zonées hors outre-mer (0,55), boîte à brevets à 17,5~\% (0,2 à 0,6) : 3,5 à 4,9 & petites sociétés ; grands groupes de recherche ; 1~065 utilisateurs de la boîte à brevets \\
Financement par le taux (28~\%) & 5,1 à 8,7 ; perte sèche +0,04 sur le capital et 0 à +3,8 selon le coût réel du transfert de bénéfices & toutes les sociétés, surtout les profits mobiles \\
Taux réduit maintenu & il manque 2,2 à 3,1 ; la dette de transition n'est jamais remboursée & --- \\
Taux réduit réservé aux sociétés opérationnelles & rendement non chiffrable (aucune ventilation publiée) & holdings patrimoniales, sociétés immobilières à l'IS \\
Excédent une fois la dette remboursée & rendu par une baisse du taux normal de 0,6 à 1,2 point & --- \\
Contribution exceptionnelle & s'éteint (référence) : $-$7 par rapport à 2026, $-$5 par rapport au PLF 2027 & environ 450 grands groupes \\
\bottomrule
\end{tabularx}
\end{table}

\textbf{Le taux ou l'assiette.} Une hausse du taux normal de 25 à 28~\% rapporterait 8,7~Md€ à bénéfices
inchangés et 5,1~Md€ si les bénéfices déclarés baissent de 1,5~\% par point, l'ordre de grandeur trouvé sur les
filiales de multinationales (Beer, de Mooij et Liu, 2020). L'objection n'est pas son montant, mais ce qu'elle
frappe : la déduction immédiate cesse d'imposer le rendement normal des investissements nouveaux, alors que le
taux porte d'abord sur les profits exceptionnels, qui sont aussi les plus mobiles. Un groupe choisit le pays
où il installe un projet très rentable selon le taux moyen effectif, et celui où il déclare ses bénéfices
selon le taux affiché. Si les bénéfices déplacés sont une perte réelle pour la France, la perte sèche de
28~\% atteint plusieurs milliards ; si ce ne sont que des écritures, elle est presque nulle (tableau~\ref{tab:is_arb}).
Le financement par l'assiette (3,5 à 4,9~Md€) dépasse le coût permanent de la déduction ; l'excédent rembourse
la transition, puis peut être rendu par une baisse du taux normal de 0,6 à 1,2 point.

\textbf{La transition} est la vraie difficulté. Une déduction totale dès 2027, financée par l'assiette,
creuserait une dette de 17 à 45~Md€, intérêts compris, remboursée en 10 à 28~ans. Un calendrier voté dès 2027,
50~\% jusqu'en 2032 puis 100~\%, limite la dette à 2 à 32~Md€, remboursée en 4 à 25~ans ; il évite de laisser le
second palier à l'appréciation des budgets futurs, ce qui reproduirait l'instabilité reprochée à l'IS.

\textbf{La fin du taux réduit : qui y perd.} Les sociétés dont le chiffre d'affaires est inférieur à 10~M€ et
le capital détenu à 75~\% par des personnes physiques paient 15~\% au lieu de 25~\% sur les 42\,500 premiers
euros de bénéfice ; depuis une décision du Conseil d'État de 2025, le chiffre d'affaires d'une société
membre d'un groupe s'apprécie au niveau du groupe. Le dispositif coûte environ 2,7~Md€ (bénéfices de 2022).
Pour une société au-dessus du plafond, il revient à verser chaque année 4\,250~€, quoi qu'elle fasse : la
somme ne dépend d'aucune décision, sauf celle de rester petite ou de se découper (on observe un pic de
sociétés juste sous le plafond). Mais beaucoup de bénéficiaires sont sous le plafond (environ 24\,000~€ de
bénéfice imposé à 15~\% en moyenne en 2018) : pour eux, la suppression relève l'impôt de chaque euro
supplémentaire, de deux tiers en tout (3\,000~€ pour 30\,000~€ de bénéfice), et renchérit tout leur capital.
Bernini et Treibich (2016) trouvent que l'introduction du taux réduit en 2001-2003 a accru de 29~\% le
capital des petites sociétés et leurs exportations, mais ralenti leur emploi (résumé seulement lu). La
déduction immédiate ne compense qu'une partie : pendant la phase à 50~\%, elle avance 2\,300~€ d'impôt la
première année pour 20\,000~€ d'équipements, et ne compense durablement la perte qu'au-delà de 220\,000~€
d'investissement annuel (110\,000~€ à 100~\%). Perdent le plus les sociétés qui investissent peu en
équipements : professions libérales, conseil, petites holdings et sociétés immobilières à l'IS ; aucune
statistique ne ventile la base imposée à 15~\% par type de société. Pour l'associé qui se fait distribuer le
bénéfice, le taux cumulé de l'IS et du PFU sur les 42\,500 premiers euros passe de 41,7 à 48,6~\%, celui de
toutes les autres sociétés.

\textbf{Ce que la référence écarte.} Un barème progressif de l'IS : la taille d'une entreprise ne dit rien de
la richesse de ses propriétaires, et les seuils multiplieraient les découpages. Un taux minimal national :
le minimum mondial de 15~\% s'applique déjà aux groupes de plus de 750~M€ de chiffre d'affaires (loi de
finances pour 2024, art.~33) ; pour les autres, un minimum assis sur l'impôt de l'année frapperait les
entreprises qui investissent beaucoup cette année-là. La contribution exceptionnelle s'éteint : surtaxe en
partie rétroactive et renouvelée d'année en année, elle fait payer les investissements passés et rend
incertain le rendement des investissements futurs.


\section{Revenus du capital : garder le prélèvement forfaitaire unique}\label{sec:pfu}
% Partie PFU (source : ../pfu/note-pfu.tex, corrigée selon ../revue/verif_is_pfu.md)
Le prélèvement forfaitaire unique (PFU) impose depuis 2018 les dividendes, les intérêts et les plus-values
mobilières à un taux unique, 31,4~\% en 2026, au lieu du barème progressif. La référence le garde tel qu'il
est. Quand le taux a monté en 2013, les actionnaires qui contrôlent leur société ont laissé les bénéfices
dedans ; quand il a baissé en 2018, ils les ont distribués. Cette réaction porte sur le moment où l'on se
verse le revenu plus que sur l'activité, et elle dépend du taux qu'on attend pour demain : le principal
service à rendre est de ne plus le changer. La faiblesse du système n'est pas dans le taux du PFU, mais dans
le circuit des holdings.

\subsection{Le PFU en 2026}

Le PFU additionne 12,8~\% d'impôt sur le revenu et 18,6~\% de prélèvements sociaux. Depuis la LFSS pour 2026,
la CSG sur les revenus du capital est de 10,6~\%, mais elle reste à 9,2~\% pour une liste fermée de revenus
(assurance-vie, épargne logement, revenus fonciers, plus-values immobilières) : deux taux coexistent, sans
justification économique. Le contribuable peut opter pour le barème, globalement et, depuis 2026, de façon
révocable ; 2,5~millions de foyers l'ont fait sur leurs revenus de 2024. En haut de la distribution, la
contribution exceptionnelle sur les hauts revenus (CEHR) ajoute 3 puis 4~points au-delà de 250\,000~€ de
revenu (500\,000~€ pour un couple), et la contribution différentielle (CDHR) garantit un impôt d'au moins
20~\% du revenu, ce qui porte à 38,6~\% le taux d'un foyer qui vit de dividendes ; elle est prorogée jusqu'à ce
que le déficit du budget de l'État repasse sous 3~\% du PIB.

Le PFU ne s'apprécie qu'avec l'IS déjà payé par la société. Sur 100~€ de bénéfice, une société paie 25~€
d'IS (25,8~€ si son impôt dépasse 763\,000~€, avec la contribution sociale) ; si elle distribue le reste,
l'actionnaire paie 31,4~\%. Le bénéfice distribué supporte donc 48,6 à 49,1~\% au total, et jusqu'à 54,5~\%
pour un foyer à la CDHR (figure~\ref{fig:taux}), proche du taux marginal supérieur sur un salaire (55,4~\%
hors cotisations sociales). C'est élevé en Europe : la France est sixième sur 35 pour le taux personnel
supérieur sur les dividendes (Tax Foundation), et l'OCDE trouve, à cinq fois le salaire moyen, un taux intégré
de 49~\% que seuls le Danemark et l'Espagne dépassent. Si la société verse son dividende à une holding de
l'actionnaire qui le garde, la holding ne paie que 1,25~\% (régime mère-fille), et le bénéfice n'a supporté
que 26,8~\%. C'est un report, pas un taux final : distribué un jour, il paiera le PFU. Imposer le bénéfice
quand il sort des sociétés est d'ailleurs ce qui gêne le moins la croissance, puisque le capital circule
d'une entreprise à l'autre sans impôt. Le report ne devient une exonération que si la sortie n'a jamais lieu,
ce que permet l'effacement au décès.

\figbloc{pfu_01_taux_cumule}{Impôt total sur 100~€ de bénéfice, selon ce qu'en fait l'actionnaire (2026)}{fig:taux}

\subsection{Ce que 2013 et 2018 ont appris}

En 2013, l'option pour un prélèvement libératoire disparaît et tous les dividendes passent au barème. Les
foyers les plus touchés déclarent 40~\% de dividendes en moins, sans hausse de leurs salaires ; les
entreprises qu'ils détiennent distribuent moins, accumulent des actifs financiers et n'investissent pas
moins (Bach et al., 2019). Le taux marginal supérieur intégré n'était pourtant passé que de 57,6 à 60,1~\% :
les actionnaires ont jugé la hausse provisoire et attendu. Pour une société qui investit ses bénéfices,
l'impôt sur le dividende ne change pas le choix d'investir : 100~€ distribués aujourd'hui rapportent 68,6~€
nets ; réinvestis à 5~\%, ils donneront 105~€ l'an prochain, soit 72,0~€ nets, et le PFU prend la même fraction
dans les deux cas. Il change seulement la date et la forme du revenu. Le raisonnement ne vaut pas pour une
jeune entreprise qui lève des fonds propres à l'extérieur ; le seul effet réel trouvé après 2018 est une
création d'entreprises un peu plus rapide dans les secteurs exposés (IPP, 2023), et la baisse américaine de
2003 n'a pas davantage modifié l'investissement (Yagan, 2015). En 2018, les dividendes remontent (+15~\% dès
la première année) et le coût de la réforme tombe de 0,9 à 0,4~Md€ selon l'IPP (0,74~Md€, voire un gain
dans l'hypothèse haute, selon l'Insee). Mais 2013 et 2018 forment un seul épisode : les bénéfices retenus de
2013 à 2017 ont été distribués à partir de 2018, et l'« autofinancement » (France Stratégie, 2023) ne vaut que
pour les premières années, rattrapage compris.

Le taux qui maximise les recettes sur les revenus du capital, prélèvements sociaux compris, est estimé entre
43~\% et 57~\% (Lefebvre, Lehmann et Sicsic, 2025). Le PFU est en dessous, mais la marge est mince : avec
l'élasticité estimée (0,77), un passage de 31,4 à 33~\% ne garde qu'un peu moins des deux tiers de son
rendement mécanique. Cette élasticité ne mesure pas une perte sèche : la réaction consiste surtout à laisser
le bénéfice dans la société, ce qui reporte l'impôt sans détruire d'activité. 1~\% des foyers déclarent
96~\% des dividendes ; l'exil est faible (un point d'imposition effective en plus fait partir à long terme 0,02
à 0,23~\% du centile le plus riche ; Bach et al., 2025) et non mesuré au-delà. La commission d'enquête de
l'Assemblée sur les plus hauts patrimoines recommande de « maintenir la stabilité » des taux (rapport
n°~3056, 2026).

\subsection{Les arbitrages}

\textbf{Le taux.} Un PFU à 25~\% coûterait au plus 5,1~Md€ à comportements inchangés (un majorant, à cause des
foyers qui optent pour le barème) et 3,2 à 3,9~Md€ après réaction. Le gain irait d'abord au centile le plus
aisé (44~\% du gain du PFU en 2018 selon le Gouvernement). Sa perte sèche dépend de la part de la réaction qui
échappe définitivement à l'impôt plutôt que d'être différée : si elle est nulle, la baisse coûte le prix de
3,7~Md€ à lever ailleurs, environ +1,4~Md€ de perte sèche ; si toute la réaction échappait à l'impôt, elle
serait à peu près neutre. Une hausse rapporterait peu : 1,2~Md€ nets pour 3~points selon le CAE.

\textbf{Le report de l'apport-cession.} Un entrepreneur qui apporte ses titres à une holding avant de les
vendre place sa plus-value en report : elle n'est imposée que si la holding vend sans réinvestir au moins
70~\% du produit. Au décès, le report devient une exonération définitive ; en cas de donation, il s'efface si
le donataire ne vend pas pendant six ans. 33,6~Md€ de plus-values ont été placées en report en 2023 et
22,2~Md€ en 2024, plus que les plus-values imposées au PFU (18,8~Md€) ; le stock atteint 153~Md€, les deux tiers
des entrées de 2023 se concentrent sur 84 foyers, et seuls 0,4 à 1,5~Md€ de reports expirent chaque année
(rapport n°~3056). Trois traitements sont possibles :
\begin{itemize}\setlength{\itemsep}{1pt}
\item l'expiration à toute transmission (PLF 2027, art.~5) : 0,5~Md€ annoncés pour 2027 et au moins 2~Md€ par
an en régime de croisière, mais elle taxe une structure plutôt qu'un gain réalisé (une holding familiale
transmise sous pacte Dutreil paierait sans avoir rien vendu ; l'étalement sur cinq ans ne vise que l'impôt
sur le revenu) et alourdit de 90 à 100~M€ chacun la facture des plus gros apporteurs ;
\item l'expiration seulement quand la holding vend ;
\item la transmission du report aux héritiers sans effacement : la plus-value est imposée quand le capital
sort réellement des sociétés, sans impôt sans liquidité ni hausse de facture à la transmission ; elle
rapporte moins et plus tard, et sa solidité juridique reste à expertiser.
\end{itemize}

\textbf{Les autres choix.} Dans l'impôt fusionné, la CEHR disparaît pour les revenus du barème ; la garder sur
les revenus du capital, sous la forme d'une majoration de 3 puis 4~points du PFU, maintient le taux actuel
(tant que la CDHR existe, un foyer qui vit de dividendes, déjà au plancher de 20~\%, n'y gagnerait de toute
façon rien). Durcir la CDHR, comme le propose la commission d'enquête, alourdirait la facture des mêmes
foyers. Aligner les deux taux de prélèvements sociaux à 18,6~\% rapporterait de l'ordre de 1,3~Md€, payés
surtout par la classe moyenne épargnante. Un prélèvement modéré sur la trésorerie passive des holdings
(l'Irlande taxe à 20~\% les revenus passifs non distribués) viserait le bon problème, mais alourdirait lui
aussi la facture des grandes fortunes. Un impôt de 25~\% sur les revenus non réinvestis des holdings se
heurte à la directive mère-fille et à la jurisprudence constitutionnelle (imposer chez l'associé des
bénéfices qu'il n'a pas perçus), et taxerait le capital avant qu'il sorte des sociétés. Enfin, la correction
de principe pour une détention longue est d'imposer le gain réel, en indexant le prix d'acquisition sur
l'inflation, plutôt que de rétablir des abattements pour durée de détention, qui avaient compliqué l'impôt
de 2013 à 2017.


\section{L'ensemble}\label{sec:ensemble}
% Synthèse PROVISOIRE
Cette section rassemblera le bilan des quatre réformes une fois la partie sur la fusion recalculée. En l'état,
pour les trois autres :

\begin{table}[!htbp]\centering\scriptsize
\caption{Bilan provisoire (Md€ par an, régime permanent)}\label{tab:ensemble}
\begin{tabularx}{\linewidth}{@{}>{\raggedright\arraybackslash}p{3.6cm}>{\raggedright\arraybackslash}p{2.6cm}>{\raggedright\arraybackslash}p{2.6cm}>{\raggedright\arraybackslash}X@{}}
\toprule
Réforme & Recette & Perte sèche & Grandes fortunes \\
\midrule
Fusion IR-CSG-CRDS & 0 par construction & en cours de recalcul & dépend de la répartition du capital entre tranches \\
Immobilier : droits $-$50~\%, deux taux, taux national sur le sol & 0 (transition : hausse temporaire de la taxe actuelle) & environ $-$2 ($-$1 à $-$5) & taxe foncière des biens chers en hausse, droits en baisse à l'achat \\
IS : déduction immédiate, financée par l'assiette & +0,6 à +1,2 point de taux à rendre une fois la dette de transition (2 à 32) remboursée & $-$0,05 à $-$0,16 (hors fin du taux réduit : 0 à +0,8) & presque rien \\
PFU maintenu, report d'apport-cession transmis sans effacement & 0 aujourd'hui, recette différée ensuite & environ 0 & pas de hausse à la transmission \\
\bottomrule
\end{tabularx}
\end{table}


\vspace{8pt}
{\scriptsize\setlength{\parskip}{2pt}\textbf{Références.}
André et Meslin (2023), \emph{Insee Analyses} n°~91.
André et Sireyjol (2021), « Imposition des couples et des familles », \emph{Économie et Statistique} 526-527.
Assemblée nationale (2026), commission d'enquête sur l'imposition des plus hauts patrimoines et des revenus les
plus élevés, rapport n°~3056.
Avenancio-León et Howard (2022), « The Assessment Gap », \emph{QJE} 137(3).
Bach, Bozio, Fabre, Guillouzouic, Leroy et Malgouyres (2019), note IPP n°~46.
Bach, Bozio, Grimprel, Guillouzouic, Landais et Malgouyres (2025), Focus du CAE n°~118.
Bachelet et Poulhès (2017), CGDD.
Banzhaf et Lavery (2010), \emph{Journal of Urban Economics} 67(2).
Beer, de Mooij et Liu (2020), \emph{Journal of Economic Surveys} 34(3).
Bérard et Trannoy (2018), \emph{Économie et Statistique} n°~500-501-502.
Bernini et Treibich (2016), \emph{Industrial and Corporate Change} 25(5).
Berry (2021), « Reassessing the Property Tax », Harris School.
Best et Kleven (2018), \emph{Review of Economic Studies} 85(1).
Bontemps, Cherbonnier et Magnac (2023), TSE WP 23-1486.
Cabannes, Houdré et Landais (2014), \emph{Économie et Statistique} 467-468.
CAE (2025, 2026), Focus n°~124 et n°~147.
Chetty (2009), « Is the Taxable Income Elasticity Sufficient to Calculate Deadweight Loss? », \emph{AEJ:
Economic Policy} 1(2).
Chetty, Friedman, Olsen et Pistaferri (2011), \emph{QJE} 126(2).
Chodorow-Reich, Smith, Zidar et Zwick (2026), NBER WP 32180.
Conseil constitutionnel, décisions 2000-437 DC, 2012-662 DC et 2019-796 DC.
Conseil des prélèvements obligatoires (2016, 2023, 2025).
Dachis, Duranton et Turner (2012), \emph{Journal of Economic Geography} 12(2).
Dolls, Fuest, Krolage et Neumeier (2025), \emph{Journal of Urban Economics} 145.
Eerola, Harjunen, Lyytikäinen et Saarimaa (2021), \emph{Journal of Urban Economics} 124.
Feldstein (1999), \emph{REStat} 81(4).
Ferreira (2010), \emph{Journal of Public Economics} 94(9-10).
France Stratégie (2023), comité d'évaluation des réformes de la fiscalité du capital, quatrième rapport.
Fritzsche et Vandrei (2019), \emph{Regional Science and Urban Economics} 74.
Hines (1999), « Three Sides of Harberger Triangles », \emph{Journal of Economic Perspectives} 13(2).
HMRC (2023), évaluation de l'allègement pour primo-accédants.
Høj, Jørgensen et Schou (2018), \emph{Fiscal Studies} 39(2).
Hourani, Millar-Powell, Perret et Ramm (2023), OCDE, \emph{Taxation Working Paper} n°~65.
IFS (2023), \emph{Green Budget}, ch.~10.
IPP (2023), rapport n°~47.
Kleven et Kreiner (2006), « The Marginal Cost of Public Funds », \emph{Journal of Public Economics} 90(10-11).
Koster et Pinchbeck (2022), \emph{AEJ: Economic Policy} 14(1).
Lefebvre, Lehmann et Sicsic (2025), \emph{Scandinavian Journal of Economics} (IZA DP 16112, 2023).
Lutz (2015), \emph{AEJ: Economic Policy} 7(1).
Määttänen et Terviö (2022), \emph{Economic Journal} 132(644).
Maffini, Xing et Devereux (2019), \emph{AEJ: Economic Policy} 11(3).
Martínez, Saez et Siegenthaler (2021), \emph{AER} 111(2).
Morin (2014), \emph{Insee Première} 1492.
Nielsson, Wroblewski et Yding (2024), document de travail.
OCDE, \emph{Revenue Statistics}, \emph{Corporate Tax Statistics} (2025), \emph{Taxing Wages} (2025).
Paquier et Sicsic (2021), \emph{Économie et Statistique} n°~530-31.
Sénat (1977), M.~Blin, rapport général sur le PLF pour 1978, n°~88 ; Sénat (2026), mission d'information sur les
prélèvements obligatoires des entreprises.
Sicsic (2022), \emph{Fiscal Studies}.
Tax Foundation, \emph{International Tax Competitiveness Index} 2025 et 2026 ; \emph{Dividend Tax Rates in
Europe} 2026.
Yagan (2015), \emph{AER} 105(12).
Zwick et Mahon (2017), \emph{AER} 107(1).
Textes : LF 2024 (art.~19 et 33), LF 2025, LFSS et LF 2026, PLF 2027 (art.~5).
Détail des sources et statut de lecture : \texttt{ir\_csg/}, \texttt{foncier/notes/}, \texttt{is/notes/},
\texttt{pfu/notes/}, \texttt{revue/}.\par}


\end{document}
